% ---
% §7.6 Limitaciones de las conclusiones.
% Entradas que lo respaldan: H-002, H-024, H-025, H-035, D-013, D-014, D-026, D-039, D-042, H-048, H-055.
% ---

\subsection{Limitaciones}
\label{sec:limitaciones_conclusiones}

Las conclusiones anteriores están acotadas por las siguientes condiciones:

\begin{itemize}
  % fuente: H-002, H-025
  \item \textbf{Un solo conjunto de datos, sin calibración física.} Todos los resultados provienen del
  CBIS-DDSM, sin validación externa. Sus mamografías son películas digitalizadas sin espaciado de píxel
  ni escala de intensidad; las características de primer orden, que lideran en masas, podrían reflejar
  en parte el digitalizador de origen.
  % fuente: D-013, H-035
  \item \textbf{Doce características.} Pasar de 67 a 12 características cuesta entre 2 y 4 puntos de AUC
  con clasificadores clásicos. La comparación es justa porque todas las condiciones parten del mismo
  vector, pero ese vector ya contiene menos información que el conjunto completo.
  % fuente: D-014, D-031, H-048
  \item \textbf{Dos codificaciones con una repetición.} Las conclusiones se refieren a los mapas ZZ y de
  Pauli con los términos Y y ZZ, con entrelazamiento completo y una repetición, y a una capa variacional
  de parámetros fijos. Otras codificaciones, en particular las diseñadas a partir de la estructura del
  problema, podrían comportarse de otro modo.
  \item \textbf{Simulación ideal.} No se modeló el ruido de los dispositivos. Todo lo que se afirma sobre
  la ejecución en un procesador real son proyecciones a partir de un muestreo ideal de las probabilidades
  exactas.
  % fuente: D-026, D-039, D-042, H-055
  \item \textbf{Alcance estadístico.} Las métricas de separabilidad del análisis principal se calcularon
  sobre 200 lesiones de entrenamiento por subconjunto; la robustez con otras submuestras y con todo el
  entrenamiento confirma la conclusión, pero la significancia de los efectos más pequeños y el orden entre
  C4 y C5 dependen de la muestra. Los valores $p$ son indicativos, porque se prueban muchas estadísticas;
  los pliegues de validación cruzada se fijaron una sola vez sobre el entrenamiento, y la correlación de
  Spearman, con cinco condiciones, es descriptiva.
\end{itemize}
